% Chapter 17: Extensions Beyond Binary Classification
%   ds_dimension — Profile C (Revelation): The 30-year gap from Natarajan to DS
%   multiclass_learning — Profile G (Framework Bridge): Multiple paradigms tried
%   agnostic_learning — Profile B-light: The agnostic fundamental theorem
%   proper_improper_separation — Profile D (Witness Construction)
% Theme: GRAMMAR GROWTH — each extension requires new primitives.
% The extends_grammar edge type is the structural backbone.

\chapter{Extensions Beyond Binary Classification}
\label{ch:extensions}

The PAC theory of \Cref{ch:pac} and its characterization through VC dimension
apply to binary classification: the label set is $\{0,1\}$, the hypothesis
class is $\mathcal{H} \subseteq \{0,1\}^X$, and the loss is $0$-$1$.  This is
the setting where the theory is cleanest and most complete.  But most learning
problems are not binary classification.  Labels may come from a finite set of
$k > 2$ classes; the target may be a real-valued function; the data may be
noisy; the realizability assumption may fail.

This chapter asks: how does the binary theory extend to these richer settings?
The answer is not ``straightforwardly.''  Each extension requires
\emph{new mathematical primitives}: new notions of shattering, new complexity
measures, new assumptions.  The binary theory generalizes by \emph{grammar
growth}: substituting richer labels into the old definitions is never enough,
because each setting forces a primitive the old language cannot name.  In the
concept graph,
this is captured by edges of type \rel{extends\_grammar}: the source dimension
extends the grammar of the VC dimension by introducing primitives that have no
analogue in the binary case.

\medskip

The chapter has six sections.

\begin{enumerate}[leftmargin=*,itemsep=2pt]
\item \textbf{Multiclass learning} (\Cref{sec:multiclass-ext}): the Natarajan
  dimension, the DS dimension, and the resolution of a thirty-year open problem.
\item \textbf{Real-valued functions} (\Cref{sec:real-valued}):
  pseudodimension and fat-shattering dimension.
\item \textbf{Agnostic learning} (\Cref{sec:agnostic-ext}): dropping realizability
  and the agnostic fundamental theorem.
\item \textbf{Noise-tolerant and partial concept learning}
  (\Cref{sec:noise-partial}): classification noise, malicious noise,
  and promise problems.
\item \textbf{Proper versus improper learning}
  (\Cref{sec:proper-improper}): the computational gap.
\item \textbf{The \rel{extends\_grammar} pattern}
  (\Cref{sec:extends-grammar}): what the extensions have in common.
\end{enumerate}


% ================================================================
%     CONCEPT MAP FIGURE
% ================================================================

\begin{figure}[ht]
\centering
\begin{tikzpicture}[
  >={Stealth[length=5pt]},
  concept/.style={
    rectangle, rounded corners=3pt, draw, thick,
    fill=layergray, font=\small\sffamily,
    minimum width=2.3cm, minimum height=0.7cm,
    align=center
  },
  % COLOR = relationship type ; LINE = evidence grade (the second axis)
  chr_solid/.style  ={->, defblue,    thick},            % characterizes, kernel-verified
  chr_lit/.style    ={->, defblue,    thick, dashed},    % characterizes, literature
  ext_concept/.style={->, edgeorange, thick, dotted},    % extends_grammar, conceptual
  ss_lit/.style     ={->, thmred,     thick, dashed},    % strictly_stronger, literature
  res_concept/.style={->, sepgreen,   thick, dotted},    % restricts, conceptual
  lb_lit/.style     ={->, boundbrown, thick, dashed},    % lower_bounds, literature
]

% Layout by structure: two hubs (vc, pac) on a horizontal spine; the multiclass
% cycle (nat, ds) fans UP, the real-valued branch (pdim, fat) fans DOWN, the
% settings (agn, pi, noise) sit on the right. Planar straight-line drawing =
% ZERO crossings. Two codes: color = relationship type, line = evidence grade.

% --- hubs ---
\node[concept] (vc)    at (0,0)      {VC\\Dimension};
\node[concept] (pac)   at (6.5,0)    {PAC\\Learning};
% --- multiclass branch (fans up) ---
\node[concept] (nat)   at (2.2,1.9)  {Natarajan\\Dimension};
\node[concept] (ds)    at (4.3,3.1)  {DS\\Dimension};
% --- real-valued branch (fans down-left) ---
\node[concept] (pdim)  at (1.8,-1.9) {Pseudo-\\dimension};
\node[concept] (fat)   at (2.8,-3.1) {Fat-Shattering\\Dimension};
% --- agnostic: small triangle just under the verified spine (clears fat->pac) ---
\node[concept] (agn)   at (3.5,-0.8) {Agnostic\\PAC};
% --- settings (right) ---
\node[concept] (noise) at (8.7,1.4)  {Noise-Tolerant\\Learning};
\node[concept] (pi)    at (8.7,-1.5) {Proper vs\\Improper};

% --- edges: color = relationship type, line = evidence grade ---
\draw[chr_solid]   (vc)    -- (pac);   % VERIFIED rung: fundamental_theorem (sc=0)
\draw[ext_concept] (nat)   -- (vc);
\draw[ext_concept] (pdim)  -- (vc);
\draw[ext_concept] (noise) -- (pac);   % NEW edge: connects the formerly floating noise node
\draw[ss_lit]      (ds)    -- (nat);
\draw[chr_lit]     (ds)    -- (pac);
\draw[chr_lit]     (fat)   -- (pac);
\draw[chr_lit]     (vc)    -- (agn);
\draw[res_concept] (agn)   -- (pac);
\draw[lb_lit]      (pi)    -- (pac);

\end{tikzpicture}

\caption[Concept map: binary PAC generalizes by grammar growth]{Concept map for
the chapter: color marks the relationship type and line style the evidence grade.
Binary PAC theory generalizes by grammar growth, each extension adding a
shattering primitive with no binary analogue, so substituting richer labels into
the binary definitions never suffices.  The one solid edge is VC dimension
characterizing binary PAC learning (\leanname{fundamental_theorem}); the
multiclass, real-valued, agnostic, and proper-versus-improper characterizations
are literature results.}
\label{fig:ch17-concept-map}
\end{figure}


% ================================================================
% SECTION 1: MULTICLASS LEARNING
% ================================================================

\section{Multiclass Learning: From Natarajan to DS}
\label{sec:multiclass-ext}

Let $Y = [k] = \{1, 2, \ldots, k\}$ for some $k \geq 2$, and let
$\mathcal{H} \subseteq Y^X$ be a multiclass hypothesis class.  The question
is: what characterizes PAC learnability of $\mathcal{H}$?

For $k = 2$, the answer is the VC dimension.  For $k > 2$, the answer was only
fully resolved in 2022, after a thirty-year journey through two competing
dimensions, each introducing a new shattering primitive.  That journey, and its
resolution, is the narrative backbone of this section.


\subsection{The Natarajan Dimension}

Natarajan~\cite{Natarajan1989} introduced the first multiclass complexity
measure by generalizing shattering to require two functions instead of one.

\begin{defn}[Natarajan Dimension {\cite{Natarajan1989}}]
\label{def:natarajan-dim}
The \emph{Natarajan dimension} $\Ndim(\mathcal{H})$ of a multiclass hypothesis
class $\mathcal{H} \subseteq Y^X$ is the largest size $d$ of a set
$\{x_1, \ldots, x_d\} \subseteq X$ for which there exist two functions
$f_0, f_1 \colon \{x_1, \ldots, x_d\} \to Y$ satisfying:
\begin{enumerate}[leftmargin=*,nosep]
\item $f_0(x_i) \neq f_1(x_i)$ for all $i \in [d]$, and
\item for every binary string $b \in \{0,1\}^d$, there exists $h \in
  \mathcal{H}$ with $h(x_i) = f_{b_i}(x_i)$ for all $i$.\formal{FLT_Proofs.Complexity.Structures.NatarajanDim}
\end{enumerate}
\end{defn}

The new primitive is the \emph{pair of functions} $(f_0, f_1)$.  In binary
classification, $f_0$ and $f_1$ are forced to be the constant functions $0$
and $1$, and the definition reduces to VC shattering.  In the multiclass case,
the choice of which two labels to use at each point is part of the
combinatorial structure.

\begin{example}[Natarajan shattering with $k=3$]
\label{ex:natarajan-k3}
Let $X = \{a, b\}$, $Y = \{1,2,3\}$.  Define $f_0(a)=1$, $f_0(b)=2$
and $f_1(a)=3$, $f_1(b)=1$.  To N-shatter $\{a,b\}$, we need four
hypotheses realizing every binary string:
\[
  h_{00}(a)=1,\, h_{00}(b)=2;\quad
  h_{01}(a)=1,\, h_{01}(b)=1;\quad
  h_{10}(a)=3,\, h_{10}(b)=2;\quad
  h_{11}(a)=3,\, h_{11}(b)=1.
\]
Any $\mathcal{H} \supseteq \{h_{00}, h_{01}, h_{10}, h_{11}\}$ has $\Ndim \geq 2$.
\end{example}

\begin{historical}
\textbf{The Natarajan gap (1989--2022).}
Natarajan showed that $\Ndim(\mathcal{H}) < \infty$ is \emph{sufficient}
for PAC learnability and provided sample complexity bounds.  The best
known bounds, due to Ben-David et al.~\cite{BenDavidCesaBianchiHausslerLong1995},
were:
\[
  C_1 \cdot \frac{\Ndim}{\varepsilon}
  \;\leq\; m_{\mathrm{PAC}}(\varepsilon,\delta)
  \;\leq\; C_2 \cdot \frac{\Ndim \cdot \ln k \cdot \ln(1/\varepsilon)}{\varepsilon}.
\]
Two gaps persisted.  First, the $\Theta(\log k)$ factor between upper
and lower bounds.  Second, and more fundamentally: is finiteness of $\Ndim$
also \emph{necessary}?  For finite label sets $|Y| = k < \infty$, the answer
is yes (the $\log k$ factor is at most a polynomial overhead).  But when
$|Y| = \infty$, a class can have $\Ndim = 1$ and yet fail to be PAC learnable.
This was the open problem that persisted for over thirty years.
\end{historical}


The \nde{natarajan\_dimension} extends the grammar of the \nde{vc\_dimension} by
the two-function shattering primitive $(f_0, f_1)$, and leaves open a tight
multiclass sample complexity without the $\log k$ factor.


\subsection{The DS Dimension: A Thirty-Year Resolution}
\label{subsec:ds-dim}

The correct characterization of multiclass PAC learnability was established
by Brukhim et al.~\cite{BrukhimCarmonDinurMoranYehudayoff2022} using the
\emph{DS dimension}, named after Daniely and
Shalev-Shwartz~\cite{DanielyShalevShwartz2014} who introduced the
underlying concept of pseudo-cubes in one-inclusion hypergraphs.

\begin{defn}[One-Inclusion Hypergraph]
\label{def:oig}
Let $\mathcal{H} \subseteq Y^X$ and let $S = \{x_1, \ldots, x_n\} \subseteq X$
be a finite set.  The \emph{one-inclusion hypergraph} $\mathrm{OIG}(\mathcal{H}, S)$
has vertex set $\mathcal{H}|_S = \{(h(x_1), \ldots, h(x_n)) : h \in \mathcal{H}\}$
and, for each coordinate $i \in [n]$, a hyperedge connecting all vertices that
agree on all coordinates except possibly coordinate $i$.
\end{defn}

\begin{defn}[DS Dimension {\cite{DanielyShalevShwartz2014,BrukhimCarmonDinurMoranYehudayoff2022}}]
\label{def:ds-dim-ext}
A set $\{x_1, \ldots, x_d\} \subseteq X$ is \emph{DS-shattered} by
$\mathcal{H} \subseteq Y^X$ if there exists a function
$f \colon \{x_1, \ldots, x_d\} \to Y$ such that for every $i \in [d]$,
there exists $h_i \in \mathcal{H}$ with:
\begin{enumerate}[leftmargin=*,nosep]
\item $h_i(x_i) \neq f(x_i)$, and
\item $h_i(x_j) = f(x_j)$ for all $j \neq i$.
\end{enumerate}
The \emph{DS dimension} $\DSdim(\mathcal{H})$ is the size of the largest
DS-shattered set.
\end{defn}

The structural content of DS-shattering is different from Natarajan shattering
in a revealing way.  Natarajan shattering asks for a \emph{global} pair of
labeling functions with \emph{all} $2^d$ interpolations realized.
DS-shattering asks for a single ``default'' labeling $f$ and, for each
point, a \emph{local} witness that deviates at exactly that point.  The
requirement is weaker per point (only one deviation, not arbitrary binary
combinations) but more structural: it encodes a property of the
one-inclusion hypergraph of $\mathcal{H}$.

\begin{figure}[ht]
\centering
\begin{tikzpicture}[
  scale=0.82,
  point/.style={circle, fill=defblue, inner sep=1.8pt},
  hyp/.style={font=\small},
  agree/.style={circle, fill=defblue!20, draw=defblue, inner sep=1.8pt},
  deviate/.style={circle, fill=thmred!20, draw=thmred, inner sep=1.8pt},
  brace/.style={decorate, decoration={brace, amplitude=5pt}},
  oigv/.style={circle, draw=defblue, fill=defblue!12, inner sep=2.4pt, font=\scriptsize},
  cedge/.style={defblue, thick}
]
% ===== LEFT PANEL: the DS witness table (the definition, d = 3) =====
\node[font=\bfseries\small] at (1.5, 3.15) {Witness table};

\node[hyp] at (-1.5, 2.0) {Default $f$:};
\foreach \i/\lbl in {1/a, 2/b, 3/c} {
  \node[point] (x\i) at (1.5*\i, 2.0) {};
  \node[above=3pt] at (x\i) {$x_{\i}$};
  \node[below=5pt, font=\small, text=defblue] at (x\i) {$\lbl$};
}

% Witness h1 : deviates at x1 only
\node[hyp] at (-1.5, 0.6) {$h_1$:};
\node[deviate] (h11) at (1.5, 0.6) {};  \node[below=5pt, font=\small, text=thmred] at (h11) {$\neq a$};
\node[agree]   (h12) at (3.0, 0.6) {};  \node[below=5pt, font=\small, text=defblue] at (h12) {$b$};
\node[agree]   (h13) at (4.5, 0.6) {};  \node[below=5pt, font=\small, text=defblue] at (h13) {$c$};

% Witness h2 : deviates at x2 only
\node[hyp] at (-1.5, -0.4) {$h_2$:};
\node[agree]   (h21) at (1.5, -0.4) {}; \node[below=5pt, font=\small, text=defblue] at (h21) {$a$};
\node[deviate] (h22) at (3.0, -0.4) {}; \node[below=5pt, font=\small, text=thmred] at (h22) {$\neq b$};
\node[agree]   (h23) at (4.5, -0.4) {}; \node[below=5pt, font=\small, text=defblue] at (h23) {$c$};

% Witness h3 : deviates at x3 only
\node[hyp] at (-1.5, -1.4) {$h_3$:};
\node[agree]   (h31) at (1.5, -1.4) {}; \node[below=5pt, font=\small, text=defblue] at (h31) {$a$};
\node[agree]   (h32) at (3.0, -1.4) {}; \node[below=5pt, font=\small, text=defblue] at (h32) {$b$};
\node[deviate] (h33) at (4.5, -1.4) {}; \node[below=5pt, font=\small, text=thmred] at (h33) {$\neq c$};

\draw[brace] (5.15, 0.9) -- (5.15, -1.7)
  node[midway, right=6pt, font=\scriptsize, align=left]
  {Each $h_i$ deviates\\from $f$ at $x_i$ only};

% ===== separator =====
\draw[gray!45, thin] (7.05, 3.3) -- (7.05, -2.5);

% ===== RIGHT PANEL: the one-inclusion hypergraph (star = the witness) =====
% OIG vertices = realized label-vectors; coordinate-i hyperedge joins vectors
% agreeing off coord i. Here f and h_i agree off coord i, so the edges are
% exactly {f,h_i}: a star. (h_i, h_j differ at TWO coords -> not adjacent.)
\node[font=\bfseries\small] at (10.0, 3.15) {One-inclusion hypergraph};

\node[oigv, fill=defblue!24] (Of)  at (10.0, 0.55) {$f$};
\node[oigv] (Oh1) at (8.55, 1.75) {$h_1$};
\node[oigv] (Oh2) at (11.45, 1.75) {$h_2$};
\node[oigv] (Oh3) at (10.0, -1.30) {$h_3$};

\draw[cedge] (Of) -- (Oh1) node[midway, above left=-2pt, font=\scriptsize, text=defblue] {$i{=}1$};
\draw[cedge] (Of) -- (Oh2) node[midway, above right=-2pt, font=\scriptsize, text=defblue] {$i{=}2$};
\draw[cedge] (Of) -- (Oh3) node[midway, right=1pt, font=\scriptsize, text=defblue] {$i{=}3$};

\node[align=center, font=\scriptsize] at (10.0, -2.15)
  {$f$ joins each $h_i$ along edge $i$; a minimum-outdegree\\
   proper orientation (the learner) runs on this structure};

\end{tikzpicture}
\caption[DS-shattering and its star in the one-inclusion hypergraph]{%
DS-shattering of a $d$-element set (here $d = 3$), drawn twice.  On the left a
default labeling $f$ assigns $a, b, c$ to $x_1, x_2, x_3$, and each witness $h_i$
agrees with $f$ except at coordinate $i$; on the right the same data is the star
joining $f$ to each $h_i$ along edge $i$ in the one-inclusion hypergraph
(\Cref{def:oig}).  Finiteness of the largest such set, the DS dimension,
characterizes multiclass PAC learnability (\Cref{thm:multiclass-char}, Brukhim et
al.\ 2022), where Natarajan shattering does not when the label set is infinite.}
\label{fig:ds-shattering}
\end{figure}


\begin{thm}[Multiclass Characterization {\cite{BrukhimCarmonDinurMoranYehudayoff2022}}]
\label{thm:multiclass-char}
A multiclass hypothesis class $\mathcal{H} \subseteq Y^X$ is PAC learnable if
and only if $\DSdim(\mathcal{H}) < \infty$.
\end{thm}

\begin{proof}[Proof sketch]
We outline both directions.

\medskip\noindent
\textbf{Sufficiency} ($\DSdim(\mathcal{H}) < \infty \Rightarrow$ learnable).
The proof proceeds through the one-inclusion hypergraph.  For a finite
restriction $\mathcal{H}|_S$ to a sample $S$ of size $m$, the one-inclusion
hypergraph $\mathrm{OIG}(\mathcal{H}, S)$ has vertices $\mathcal{H}|_S$ and
hyperedges indexed by coordinates.  The key insight of Daniely and
Shalev-Shwartz is that a \emph{proper orientation} of this hypergraph, an
assignment of each hyperedge to one of its vertices, yields a learning
algorithm.  Specifically:

\begin{enumerate}[leftmargin=*,nosep]
\item Given a sample $S = \{(x_1, y_1), \ldots, (x_m, y_m)\}$, restrict
  $\mathcal{H}$ to $\{x_1, \ldots, x_m\}$ and build $\mathrm{OIG}(\mathcal{H}, S)$.
\item Find a proper orientation with \emph{minimum outdegree}.
\item Return the hypothesis whose vertex has minimum outdegree in the
  orientation.
\end{enumerate}

If $\DSdim(\mathcal{H}) = d < \infty$, then the one-inclusion hypergraph
of any restriction has no $d$-dimensional pseudo-cube, and a combinatorial
argument (via the Helly property of CAT(0) cube complexes) shows that
a proper orientation with outdegree at most $O(d)$ exists.  The learner's
error is bounded by $O(d/m)$, yielding PAC learnability with sample
complexity $O(d/\varepsilon)$.

\medskip\noindent
\textbf{Necessity} ($\DSdim(\mathcal{H}) = \infty \Rightarrow$ not learnable).
If $\DSdim(\mathcal{H}) = \infty$, then for every $d$, there exists a
finite set that is DS-shattered.  The proof constructs, for each $d$, a
distribution $D_d$ over $X \times Y$ such that:
\begin{itemize}[leftmargin=*,nosep]
\item The Bayes-optimal classifier has zero error.
\item Any learner given $m < d$ samples has expected error
  $\Omega(d/m)$ on $D_d$.
\end{itemize}
The distribution is supported on the DS-shattered set and uses the default
labeling $f$ as the target.  The witness hypotheses $h_1, \ldots, h_d$ act
as ``confusers'': a learner that has not seen point $x_i$ cannot distinguish
$f$ from $h_i$, because they agree everywhere except at $x_i$.  Since this
works for arbitrarily large $d$, no finite sample suffices for
$\varepsilon$-accurate learning uniformly over $D_d$.
\end{proof}


\begin{obstbox}
\textbf{Natarajan $\not\Rightarrow$ PAC learnability when $|Y| = \infty$.}

\medskip\noindent
\emph{Witness (Brukhim et al.~2022).}
There exists a hypothesis class $\mathcal{H} \subseteq Y^X$ with
$\Ndim(\mathcal{H}) = 1$ but $\DSdim(\mathcal{H}) = \infty$.  The
construction uses a hyperbolic pseudo-manifold: the domain $X$ is the
vertex set of a high-dimensional simplicial complex with hyperbolic
geometry, and $\mathcal{H}$ encodes colorings of simplices.  The
hyperbolic geometry ensures that any two points can be N-shattered
(only two labels are needed locally), but the global structure admits
arbitrarily large DS-shattered sets because the one-inclusion hypergraph
contains arbitrarily large pseudo-cubes.

\medskip\noindent
\emph{Structural lesson.}
The Natarajan dimension is a \emph{local} measure: it examines pairs of
labels at each point.  The DS dimension is a \emph{global} measure: it
examines the connectivity structure of the entire one-inclusion hypergraph.
When $|Y| = \infty$, local two-label constraints do not capture global
learnability.

\medskip\noindent
\emph{Graph edge.}  Finiteness of the \nde{natarajan\_dimension} does not imply
\nde{pac\_learning}, witnessed by the Brukhim et al.\ (2022) non-learnable class
with $\Ndim = 1$.
\end{obstbox}


The relationship between the two dimensions is:

\begin{prop}[DS vs.\ Natarajan]
\label{prop:ds-vs-nat}
For any $\mathcal{H} \subseteq Y^X$:
\[
  \Ndim(\mathcal{H}) \;\leq\; \DSdim(\mathcal{H}).
\]
Moreover, when $|Y| = k < \infty$:
\[
  \DSdim(\mathcal{H}) \;\leq\; O\!\bigl(\Ndim(\mathcal{H}) \cdot \log k\bigr).
\]
\end{prop}

\begin{proof}
For the first inequality, we show that every N-shattered set is also
DS-shattered.  Suppose $S = \{x_1, \ldots, x_d\}$ is N-shattered by a pair
$(f_0, f_1)$.  Take the default labeling $f = f_0$.  For each $i$, the binary
string $e_i$ (with a $1$ in coordinate $i$ and $0$ elsewhere) is realized by
some $h_i \in \mathcal{H}$ with $h_i(x_i) = f_1(x_i) \neq f_0(x_i) = f(x_i)$
and $h_i(x_j) = f_0(x_j) = f(x_j)$ for $j \neq i$.  Thus $f$ together with
$h_1, \ldots, h_d$ witnesses that $S$ is DS-shattered.  Every N-shattered set
is therefore DS-shattered, so $\Ndim \leq \DSdim$.  DS-shattering is the
\emph{weaker} requirement (only $d$ single-coordinate deviations, not all
$2^d$ interpolations), which is exactly why $\DSdim$ can be infinite while
$\Ndim$ stays finite, as in the obstruction above.

For the second inequality (finite $k$), suppose $\DSdim = D$.  Any
DS-shattered set provides $D$ witness hypotheses deviating from a default.
Since each deviation chooses one of $k-1$ alternative labels, a pigeonhole
argument over the $k^D$ possible deviation patterns shows
$D \leq O(\Ndim \cdot \log k)$.
\end{proof}


The \nde{ds\_dimension} characterizes \nde{pac\_learning} in the multiclass
setting and is strictly stronger than the \nde{natarajan\_dimension}, the measure
that resolved a thirty-year open problem.


% ================================================================
% SECTION 2: REAL-VALUED FUNCTIONS
% ================================================================

\section{Real-Valued Functions}
\label{sec:real-valued}

Now let $\mathcal{F} \subseteq \R^X$ be a class of real-valued functions.
The learner observes samples $(x_i, y_i)$ with $y_i = f(x_i)$ (or
$y_i \approx f(x_i)$ in the noisy case) and aims to approximate $f$.
The $0$-$1$ loss is replaced by the squared loss or absolute loss, and the
question becomes: what is the analogue of VC dimension?

Two dimensions extend the binary theory.  Both introduce \emph{thresholds} as
a new primitive.


\subsection{Pseudodimension}

\begin{defn}[Pseudodimension {\cite{Pollard1984,Haussler1992}}]
\label{def:pseudodim-ext}
The \emph{pseudodimension} $\Pdim(\mathcal{F})$ of a function class
$\mathcal{F} \subseteq \R^X$ is the VC dimension of the class of subgraphs:
\[
  \Pdim(\mathcal{F}) \;=\; \VC\!\bigl(\bigl\{
    x \mapsto \ind[f(x) \geq t] : f \in \mathcal{F},\; t \in \R
  \bigr\}\bigr).
\]
Equivalently, $\Pdim(\mathcal{F}) = d$ if there exist $x_1, \ldots, x_d \in X$
and thresholds $t_1, \ldots, t_d \in \R$ such that for every
$b \in \{0,1\}^d$, there exists $f \in \mathcal{F}$ with
$f(x_i) \geq t_i \iff b_i = 1$.\formal{FLT_Proofs.Complexity.Structures.Pseudodimension}
\end{defn}

The new primitive is the \emph{threshold vector} $(t_1, \ldots, t_d)$.
In binary classification, thresholds are unnecessary: the outputs are
already $0$ or $1$.  For real-valued functions, thresholds convert the
continuous outputs into a binary shattering condition.

\begin{example}[Linear functions]
For the class of linear functions $\mathcal{F} = \{x \mapsto w^\top x : w \in \R^d\}$
on $\R^d$, we have $\Pdim(\mathcal{F}) = d$.  Any $d$ affinely independent
points $x_1, \ldots, x_d$ can be P-shattered with thresholds $t_i = 0$:
for each $b \in \{0,1\}^d$, one can find $w$ such that $\mathrm{sign}(w^\top x_i) = (-1)^{1-b_i}$.
\end{example}

\begin{prop}[Pseudodimension and covering numbers]
\label{prop:pdim-covering}
If $\Pdim(\mathcal{F}) = d < \infty$, then for every probability measure
$P$ on $X$ and every $\varepsilon > 0$:
\[
  \mathcal{N}(\varepsilon, \mathcal{F}, L_1(P))
  \;\leq\; \left(\frac{2e}{\varepsilon}\right)^d.
\]
\end{prop}

\begin{proof}[Proof sketch]
The subgraph class $\{x \mapsto \ind[f(x) \geq t] : f \in \mathcal{F},\,
t \in \R\}$ has VC dimension $\Pdim(\mathcal{F}) = d$ by definition.  Applying
the Sauer--Shelah bound and the standard covering-number estimate for VC
classes (\Cref{ch:generalization}) to this binary class, then transferring
back to $L_1(P)$ through the layer-cake identity
$f = \int_0^1 \ind[f \geq t]\,\mathrm{d}t$, yields the stated bound.
\end{proof}

As a complexity measure, the \nde{pseudodimension} extends the grammar of the
\nde{vc\_dimension}, and the extension is the threshold primitive.


\subsection{Fat-Shattering Dimension}

Pseudodimension controls learnability in the realizable case.  For the
agnostic setting with real-valued targets, a scale-sensitive
notion is needed.

\begin{defn}[Fat-Shattering Dimension {\cite{AlonBenDavidCesaBianchiHaussler1997}}]
\label{def:fat-shattering-ext}
For $\gamma > 0$, the \emph{$\gamma$-fat-shattering dimension}
$\fatshat_\gamma(\mathcal{F})$ is the largest $d$ such that there exist
$x_1, \ldots, x_d \in X$ and thresholds $t_1, \ldots, t_d \in \R$ satisfying:
for every $b \in \{0,1\}^d$, there exists $f \in \mathcal{F}$ with
\[
  f(x_i) \geq t_i + \gamma \;\text{if } b_i = 1,
  \qquad
  f(x_i) \leq t_i - \gamma \;\text{if } b_i = 0.
\]
\end{defn}

The new primitive beyond pseudodimension is the \emph{margin}
$\gamma$.\formal{FLT_Proofs.Complexity.Structures.FatShatteringDim}
As $\gamma \to 0$, the fat-shattering dimension approaches the
pseudodimension.

\begin{figure}[ht]
\centering
\begin{tikzpicture}[scale=0.9]
  % single margin value: forbidden band is [t_i - \g, t_i + \g], width 2\g, centred on t_i
  \def\g{0.5}

  % --- Axes ---
  \draw[->] (0,-0.5) -- (7,-0.5) node[right] {$x$};
  \draw[->] (0,-0.5) -- (0,4.2) node[above] {$f(x)$};

  % --- Per point: symmetric allowed zones, threshold, x-tick + label ---
  \foreach \i/\xpos/\tval in {1/1.5/1.5, 2/3.5/2.5, 3/5.5/1.8} {
    \fill[edgeorange!22] (\xpos-0.32, \tval+\g) rectangle (\xpos+0.32, \tval+\g+0.85);
    \fill[defblue!18]    (\xpos-0.32, \tval-\g-0.85) rectangle (\xpos+0.32, \tval-\g);
    \draw[dashed, notegray] (\xpos-0.5, \tval) -- (\xpos+0.5, \tval)
      node[right, font=\tiny, notegray] {$t_{\i}$};
    \draw (\xpos, -0.55) -- (\xpos, -0.45);
    \node[font=\tiny] at (\xpos, -0.75) {$x_{\i}$};
  }

  % --- zone labels + 2gamma brace, marked once at x_1 to keep it clean ---
  \node[font=\tiny, edgeorange!80!black] at (1.5, 1.5+\g+0.45) {$\geq t_i{+}\gamma$};
  \node[font=\tiny, defblue]             at (1.5, 1.5-\g-0.45) {$\leq t_i{-}\gamma$};
  \draw[<->, thick, sepgreen] (1.5+0.42, 1.5-\g) -- (1.5+0.42, 1.5+\g)
    node[midway, right, font=\tiny, text=sepgreen] {$2\gamma$};

  % ===== FUNCTION CURVES (the witnesses that exhibit the definition) =====
  % f_A realizes b=(1,0,1): >= t_1+g, <= t_2-g, >= t_3+g
  \draw[thmred, very thick]
    plot[smooth, tension=0.7] coordinates
    {(0.3,1.9) (1.5,2.4) (2.5,2.3) (3.5,1.7) (4.5,1.9) (5.5,2.6) (6.4,2.7)};
  \fill[thmred] (1.5,2.4) circle (2.2pt);
  \fill[thmred] (3.5,1.7) circle (2.2pt);
  \fill[thmred] (5.5,2.6) circle (2.2pt);
  \node[thmred, font=\scriptsize] at (2.5,3.2) {$f_A\colon b{=}101$};

  % f_B realizes b=(0,1,0): <= t_1-g, >= t_2+g, <= t_3-g
  \draw[analogypurple, very thick, dashed]
    plot[smooth, tension=0.7] coordinates
    {(0.3,1.1) (1.5,0.6) (2.5,1.6) (3.5,3.3) (4.5,1.9) (5.5,0.9) (6.4,0.7)};
  \fill[analogypurple] (1.5,0.6) circle (2.2pt);
  \fill[analogypurple] (3.5,3.3) circle (2.2pt);
  \fill[analogypurple] (5.5,0.9) circle (2.2pt);
  \node[analogypurple, font=\scriptsize] at (4.6,3.6) {$f_B\colon b{=}010$};
\end{tikzpicture}
\caption[$\gamma$-fat-shattering: curves realizing labelings against shared
thresholds]{$\gamma$-fat-shattering of $d=3$ points: the set is
$\gamma$-fat-shattered when thresholds $t_1,t_2,t_3$ admit, for every labeling, a
function clearing the margin (at least $t_i+\gamma$ where the label is $1$, at most
$t_i-\gamma$ where it is $0$).  The forbidden band of width $2\gamma$ makes the
notion scale-sensitive; the two curves realize $101$ and $010$ against a shared
threshold vector.  Finiteness for every $\gamma$ (\leanname{FatShatteringDim})
characterizes agnostic real-valued PAC learnability (Alon et al.\ 1997).}
\label{fig:fat-shattering}
\end{figure}

\begin{thm}[Real-Valued Characterization {\cite{AlonBenDavidCesaBianchiHaussler1997}}]
\label{thm:real-valued-char}
A class $\mathcal{F} \subseteq [0,1]^X$ of bounded real-valued functions is
agnostically PAC learnable with respect to the absolute loss if and only if
$\fatshat_\gamma(\mathcal{F}) < \infty$ for all $\gamma > 0$.
\end{thm}

\begin{proof}[Proof sketch]
\textbf{Sufficiency.}  If $\fatshat_\gamma(\mathcal{F}) = d_\gamma < \infty$
for all $\gamma > 0$, then the $L_2$ covering numbers satisfy
$\log \mathcal{N}_2(\varepsilon, \mathcal{F}, x_1^n) \leq
O(d_{\varepsilon/4} \cdot \log^2(n/d_{\varepsilon/4}))$.
Finite covering numbers imply uniform convergence of empirical risk to
true risk, which yields agnostic PAC learnability.

\textbf{Necessity.}  If $\fatshat_\gamma(\mathcal{F}) = \infty$ for some
$\gamma > 0$, then for every $n$, one can $\gamma$-fat-shatter $n$ points.
By a probabilistic argument, any empirical risk minimizer on a sample of
size $m$ has expected error $\Omega(\gamma)$ on a suitably chosen
distribution, regardless of $m$.  Hence $\mathcal{F}$ is not learnable
at accuracy $\gamma$.
\end{proof}

\begin{remark}[Scale families versus single numbers]
\label{rmk:scale-family}
The fat-shattering characterization is a \emph{family} of finiteness
requirements, one for each $\gamma > 0$, where a single threshold-free number
would have sufficed in the binary case.  This is
another instance of grammar growth: the real-valued setting requires a
scale parameter that has no analogue in the binary theory.  In practice,
$\fatshat_\gamma$ is often a decreasing function of $\gamma$; larger
margins are harder to maintain, so fewer points can be fat-shattered.
\end{remark}


% ================================================================
% SECTION 3: AGNOSTIC LEARNING
% ================================================================

\section{Agnostic Learning: Dropping Realizability}
\label{sec:agnostic-ext}

All results so far have assumed \emph{realizability}: there exists some
$h^* \in \mathcal{H}$ with $\risk_D(h^*) = 0$.  In practice, models are
always wrong; the question is how wrong.  The agnostic setting, introduced
by Haussler~\cite{Haussler1992}, removes the realizability assumption entirely.

\begin{defn}[Agnostic PAC Learning {\cite{Haussler1992}}]
\label{def:agnostic-pac-ext}
A hypothesis class $\mathcal{H}$ is \emph{agnostically PAC learnable}
if there exists a learner $A$ and a function $m \colon (0,1)^2 \to \N$ such
that for every distribution $D$ over $X \times \{0,1\}$ (with \emph{no}
assumption that the Bayes-optimal classifier lies in $\mathcal{H}$), for
every $\varepsilon, \delta > 0$: given $m(\varepsilon, \delta)$ i.i.d.\
samples from $D$,
\[
  \Pr_{S \sim D^m}\!\Bigl[
    \risk_D\!\bigl(A(S)\bigr)
    \;\leq\;
    \inf_{h \in \mathcal{H}} \risk_D(h) + \varepsilon
  \Bigr] \;\geq\; 1 - \delta.
\]
\end{defn}

The learner must compete with the \emph{best hypothesis in the class},
not with the Bayes-optimal classifier.  The gap
$\inf_{h \in \mathcal{H}} \risk_D(h)$ may be positive; the learner need
only close to within $\varepsilon$ of it.

\begin{thm}[Agnostic Fundamental Theorem]
\label{thm:agnostic-ft}
For binary classification with $0$-$1$ loss, a hypothesis class
$\mathcal{H} \subseteq \{0,1\}^X$ is agnostically PAC learnable if and only
if $\VC(\mathcal{H}) < \infty$.
\end{thm}

\begin{proof}[Proof sketch]
The characterizing condition is the same as in the realizable case, finite
VC dimension, but the proof mechanism and sample complexity differ.

\medskip\noindent
\textbf{Sufficiency.}  Suppose $\VC(\mathcal{H}) = d < \infty$.  The learner
runs empirical risk minimization: $A(S) = \arg\min_{h \in \mathcal{H}}
\emprisk_S(h)$.  By the uniform convergence theorem for VC classes, with
probability at least $1 - \delta$ over a sample of size
$m = O\!\bigl(\frac{d + \log(1/\delta)}{\varepsilon^2}\bigr)$,
\[
  \sup_{h \in \mathcal{H}} \bigl|\risk_D(h) - \emprisk_S(h)\bigr|
  \;\leq\; \varepsilon/2.
\]
Let $h^* = \arg\min_{h \in \mathcal{H}} \risk_D(h)$ and $\hat{h} = A(S)$.
Then:
\begin{align*}
  \risk_D(\hat{h})
  &\leq \emprisk_S(\hat{h}) + \varepsilon/2
    &\text{(uniform convergence)}\\
  &\leq \emprisk_S(h^*) + \varepsilon/2
    &\text{(ERM chose $\hat{h}$)}\\
  &\leq \risk_D(h^*) + \varepsilon/2 + \varepsilon/2
    &\text{(uniform convergence)}\\
  &= \risk_D(h^*) + \varepsilon.
\end{align*}

\medskip\noindent
\textbf{Necessity.}  If $\VC(\mathcal{H}) = \infty$, the No-Free-Lunch
argument of the realizable case applies \textit{a fortiori}: for any learner
and any sample size $m$, there exists a realizable distribution (which is
a special case of an arbitrary distribution) on which the learner fails.
\end{proof}

\begin{separation}
\textbf{Agnostic versus realizable sample complexity.}

\medskip\noindent
The agnostic and realizable settings are characterized by the \emph{same}
condition ($\VC < \infty$), but the sample complexities differ:
\[
  m_{\mathrm{realizable}}(\varepsilon, \delta)
  = \Theta\!\left(\frac{d}{\varepsilon} + \frac{\log(1/\delta)}{\varepsilon}\right),
  \qquad
  m_{\mathrm{agnostic}}(\varepsilon, \delta)
  = \Theta\!\left(\frac{d}{\varepsilon^2} + \frac{\log(1/\delta)}{\varepsilon^2}\right).
\]
The quadratic dependence on $1/\varepsilon$ in the agnostic case is inherent: a
matching lower bound forces it, so no sharper analysis can remove it.  The bound is achieved by
distributions where the optimal hypothesis has error $1/2 - \varepsilon$,
and distinguishing it from a hypothesis with error $1/2$ requires
$\Omega(1/\varepsilon^2)$ samples by a Chernoff-bound argument.

\medskip\noindent
\emph{Graph edge.}  \nde{agnostic\_pac\_learning} restricts \nde{pac\_learning}
in the graph: agnostic PAC learning \emph{generalizes} standard PAC learning by
removing the realizability constraint.
\end{separation}


% ================================================================
% SECTION 4: NOISE AND PARTIAL CONCEPTS
% ================================================================

\section{Noise-Tolerant and Partial Concept Learning}
\label{sec:noise-partial}

\subsection{Classification Noise}

The \emph{classification noise} model (Angluin and
Laird~\cite{AngluinLaird1988}) modifies PAC learning by flipping each
label independently with probability $\eta < 1/2$: the learner observes
$(x, y \oplus z)$ where $z \sim \mathrm{Bernoulli}(\eta)$ and $y$ is the true
label.  The question is whether finite VC dimension still suffices for
learning.

\begin{thm}[Noise-tolerant PAC learning]
\label{thm:noise-tolerant}
If $\VC(\mathcal{H}) = d < \infty$ and the noise rate $\eta < 1/2$ is known,
then $\mathcal{H}$ is PAC learnable under classification noise with sample
complexity
\[
  m = O\!\left(\frac{d}{(1-2\eta)^2 \varepsilon^2}
  + \frac{\log(1/\delta)}{(1-2\eta)^2 \varepsilon^2}\right).
\]
\end{thm}

The bound is an upper bound, achieved by a noise-robust empirical risk
minimizer~\cite{AngluinLaird1988}; the $(1-2\eta)^{-2}$ dependence is also
necessary, so the rate is tight.  The sample complexity degrades as
$\eta \to 1/2$, where the noise rate approaches the information-theoretic
limit at which labels carry no signal.

\begin{remark}[SQ connection]
\label{rmk:sq-noise}
Statistical query (SQ) learnable classes are automatically noise-tolerant:
an SQ algorithm estimates expectations $\E_D[\phi(x, y)]$ to tolerance
$\tau$, and these can be simulated from noisy samples with a $1/(1-2\eta)$
overhead in sample size.  However, not all noise-tolerant classes are
SQ learnable: parities are PAC learnable (via Gaussian elimination) and
hence noise-tolerant for known $\eta$, but not SQ learnable
($\SQdim = 2^n$).
\end{remark}


\subsection{Partial Concept Learning}

In \emph{partial concept learning} (also called \emph{promise problems}),
the hypothesis class $\mathcal{H}$ is defined only on a subset
$P \subseteq X$, the \emph{promise set}.  Points outside $P$ have no
defined label, and the distribution $D$ is supported on $P$.

\begin{defn}[Partial Concept Class]
\label{def:partial-concept}
A \emph{partial concept class} is a set $\mathcal{H} \subseteq \{0, 1, *\}^X$
where $*$ denotes ``undefined.''  A distribution $D$ is \emph{compatible}
with $h \in \mathcal{H}$ if $D$ is supported on $\{x : h(x) \neq *\}$.
\end{defn}

Partial concepts arise naturally in many settings: medical diagnosis where
only symptomatic patients have definite diagnoses, or feature spaces where
certain regions are ``don't care'' zones.

Alon, Hanneke, Holzman, and Moran~\cite{AlonHannekeHolzmanMoran2021}
developed the PAC theory of partial concept classes, and the picture is
genuinely new.  A partial class is PAC learnable if and only if its
(partial) VC dimension is finite, where shattering is restricted to the
promise region; so far, the binary analogue carries over.  What does
\emph{not} carry over is the reduction to total concepts: they exhibit
partial classes of VC dimension~$1$ that admit \emph{no} total extension of
finite VC dimension, and that cannot be learned by empirical risk
minimization at all.  Partial-concept learning is therefore not a special
case of the total theory; it is a setting where proper and ERM-based
learning provably fail while improper learning succeeds, and where the
computational landscape shifts: problems easy as total concepts can become
hard under a promise constraint.


% ================================================================
% SECTION 5: PROPER VS IMPROPER
% ================================================================

\section{Proper Versus Improper Learning}
\label{sec:proper-improper}

All PAC definitions permit the learner to output \emph{any} efficiently
evaluatable hypothesis, not just one from the target class $\mathcal{H}$.
When the learner is restricted to output a hypothesis from $\mathcal{H}$
itself, we call it a \emph{proper} learner.

\begin{defn}[Proper and Improper PAC Learning]
\label{def:proper-improper-ext}
A PAC learner $A$ for class $\mathcal{H}$ is \emph{proper} if
$A(S) \in \mathcal{H}$ for all samples $S$.  It is \emph{improper} if
$A(S)$ may lie in a larger class $\mathcal{H}' \supseteq \mathcal{H}$.
\end{defn}

\begin{remark}
In the information-theoretic (sample complexity) sense, there is no
separation between proper and improper learning for standard PAC: both are
characterized by finite VC dimension, and the sample complexities are
$\Theta(d/\varepsilon)$ in both cases.  The separation is \emph{computational}.
\end{remark}

\begin{thm}[Pitt--Valiant {\cite{PittValiant1988}}]
\label{thm:pitt-valiant}
For $k \geq 2$, $k$-term DNF is not properly PAC learnable unless $\mathrm{RP} = \mathrm{NP}$.
\end{thm}

\begin{proof}[Proof]
The proof reduces graph $k$-colorability to the problem of finding a
consistent $k$-term DNF formula.

Let $G = (V, E)$ be a graph with $V = \{v_1, \ldots, v_n\}$.  We construct
a learning instance over $X = \{0,1\}^n$ as follows.

\medskip\noindent
\textbf{Positive examples.}  For each vertex $v_i$, create the example
$e_i^+ \in \{0,1\}^n$ defined by $e_i^+(j) = 1$ if $j \neq i$ and
$e_i^+(i) = 0$.  Label: positive.

\medskip\noindent
\textbf{Negative examples.}  For each edge $(v_i, v_j) \in E$, create the
example $e_{ij}^- \in \{0,1\}^n$ defined by $e_{ij}^-(l) = 1$ if
$l \notin \{i, j\}$, and $e_{ij}^-(i) = e_{ij}^-(j) = 0$.  Label: negative.

\medskip\noindent
\textbf{Claim.}  A consistent $k$-term DNF exists for this labeled sample
if and only if $G$ is $k$-colorable.

\medskip\noindent
\emph{Forward direction.}  Given a proper $k$-coloring $c \colon V \to [k]$,
define the $k$-term DNF $\phi = T_1 \vee \cdots \vee T_k$ where
$T_j = \bigwedge_{i : c(v_i) = j} \overline{x_i}$.  Each positive example
$e_i^+$ satisfies $T_{c(v_i)}$ (since $e_i^+(i) = 0$ makes the literal
$\overline{x_i}$ true, and all other literals in $T_{c(v_i)}$ are true
because $e_i^+(l) = 1$ for $l \neq i$).  Each negative example $e_{ij}^-$
falsifies every term: for any term $T_j$, since $(v_i, v_j)$ is an edge
and $c$ is proper, $c(v_i) \neq c(v_j)$, so at most one of $v_i, v_j$
is in color class $j$, and the other has $e_{ij}^-(l) = 0$ which falsifies
a literal not in $T_j$.  (A careful case analysis confirms the argument.)

\medskip\noindent
\emph{Reverse direction.}  Given a consistent $k$-term DNF
$\phi = T_1 \vee \cdots \vee T_k$, assign color $c(v_i) = \min\{j : e_i^+
\text{ satisfies } T_j\}$.  Since $e_i^+$ is positive, at least one term is
satisfied.  For any edge $(v_i, v_j)$, the negative example $e_{ij}^-$
falsifies $\phi$, which forces $c(v_i) \neq c(v_j)$ (otherwise the term
satisfied by both would also satisfy $e_{ij}^-$).  Hence $c$ is a proper
$k$-coloring.

\medskip\noindent
\textbf{Conclusion.}  A proper PAC learner for $k$-term DNF, given the
above examples (which can be generated efficiently from $G$), would find
a consistent $k$-term DNF in polynomial time, thereby solving $k$-colorability
in polynomial time.  Since $k$-colorability is NP-complete for $k \geq 3$
(and NP-hard to decide in randomized polynomial time unless
$\mathrm{RP} = \mathrm{NP}$), proper PAC learning of $k$-term DNF is
computationally hard.
\end{proof}


\begin{separation}
\textbf{Proper $\not\Rightarrow$ efficient learning: the DNF witness.}

\medskip\noindent
The same class that is hard to learn properly is easy to learn improperly.
Observe that every $k$-term DNF over $n$ variables is a $k$-CNF formula
(by distributing), and $k$-CNF is learnable in polynomial time: simply
enumerate all $O(n^k)$ candidate clauses and filter by consistency with the
data.

\medskip\noindent
\emph{Summary.} For $k = 3$: 3-term DNF is NP-hard proper, polynomial-time
improper (via 3-CNF).  This is the cleanest known computational separation
in PAC learning.

\medskip\noindent
\emph{Graph edge.}  The \nde{proper\_improper\_separation} does not imply
\nde{pac\_learning}, witnessed by $3$-term DNF: NP-hard proper, polynomial-time
improper.
\end{separation}


\begin{figure}[ht]
\centering
\begin{tikzpicture}[
  >={Stealth[length=5pt]},
  % cells: fill = learner type (improper green / proper red);
  %        border = evidence (solid = machine-checked / dashed = literature)
  cell/.style={rectangle, rounded corners=3pt, draw, thick,
    minimum width=3.5cm, minimum height=1.3cm, font=\small, align=center, inner sep=3pt},
  vcell/.style={cell},                 % solid border = kernel-verified
  lcell/.style={cell, dashed},         % dashed border = literature
  impf/.style={fill=sepgreen!12},
  propf/.style={fill=thmred!12},
  hd/.style={font=\small\bfseries\sffamily, align=center},
  rl/.style={font=\small\bfseries\sffamily, align=center, text width=2.5cm},
  verdict/.style={font=\scriptsize\itshape, text=notegray, align=center},
  rel/.style={-{Stealth[length=5pt]}, dashed, notegray, thick}
]
% --- column headers ---
\node[hd, text=sepgreen] (hi) at (0,2.05)   {Improper\\[-1pt]{\scriptsize $h \in \mathcal{H}'$}};
\node[hd, text=thmred]   (hp) at (4.7,2.05)  {Proper\\[-1pt]{\scriptsize $h \in \mathcal{H}$}};

% --- row 1: STATISTICAL layer (SOLID border = machine-checked) ---
\node[vcell, impf]  (si) at (0,0.65)   {$\Theta(d/\varepsilon)$ samples};
\node[vcell, propf] (sp) at (4.7,0.65)  {$\Theta(d/\varepsilon)$ samples};

% --- row 2: COMPUTATIONAL layer (DASHED border = literature) ---
\node[lcell, impf]  (ci) at (0,-1.35)  {Polynomial time\\{\footnotesize (via $k$-CNF)}};
\node[lcell, propf] (cp) at (4.7,-1.35) {NP-hard\\{\footnotesize (via $k$-colouring)}};

% --- row labels (left) ---
\node[rl, text=defblue]    at (-2.95,0.65)  {Statistical\\[-1pt]{\scriptsize (machine-checked)}};
\node[rl, text=boundbrown] at (-2.95,-1.35) {Computational\\[-1pt]{\scriptsize (literature)}};

% --- verdicts (right) ---
\node[verdict] at (7.3,0.65)  {same:\\no gap};
\node[verdict] at (7.3,-1.35) {differs:\\the gap};

% --- restriction relation (row 1, between columns) ---
\draw[rel] (sp.west) -- node[above, font=\scriptsize\itshape, text=notegray] {$\subseteq\mathcal{H}$} (si.east);

% --- witness (bottom) ---
\node[font=\scriptsize, align=center, text width=10cm] at (2.35,-2.75)
  {\textbf{Witness.} $3$-term DNF is one class in all four cells: the same statistically,
   yet NP-hard to learn properly and polynomial-time improperly (as $3$-CNF).};
\end{tikzpicture}
\caption[Proper versus improper learning: no statistical gap, a computational
gap]{Proper versus improper PAC learning across two layers.  In the statistical
layer (top row) there is no separation: proper and improper learning are both
characterized by finite VC dimension and both need $\Theta(d/\varepsilon)$ samples
(\leanname{fundamental_theorem}).  The gap appears only in the computational layer
(bottom row), where $3$-term DNF is NP-hard properly (Pitt and Valiant 1988,
\Cref{thm:pitt-valiant}) yet polynomial-time improperly as $3$-CNF.}
\label{fig:proper-improper}
\end{figure}


As an impossibility result, the \nde{proper\_improper\_separation} both
lower-bounds \nde{pac\_learning} and fails to imply it.  Its provenance is
Pitt--Valiant 1988, and the proof is a direct reduction from graph coloring to
consistent hypothesis finding.


% ================================================================
% SECTION 6: THE EXTENDS_GRAMMAR PATTERN
% ================================================================

\section{The \rel{extends\_grammar} Pattern}
\label{sec:extends-grammar}

The extensions of the previous sections share a common structure.  Each
begins with the VC-dimension framework and introduces new primitives that
have no analogue in binary classification:

\begin{center}\small
\begin{tabular}{llp{5.5cm}}
\toprule
\textbf{Extension} & \textbf{New Primitive} & \textbf{Why Needed} \\
\midrule
Natarajan dim. & Two-function shattering $(f_0, f_1)$
  & With $k > 2$ labels, a single bit string does not
    describe a labeling. \\[4pt]
DS dimension & One-inclusion witness
  & Global hypergraph structure governs learnability
    when $|Y| = \infty$. \\[4pt]
Pseudodimension & Threshold vector $(t_1, \ldots, t_d)$
  & Real-valued outputs must be discretized before
    shattering can be defined. \\[4pt]
Fat-shattering dim. & Scale parameter $\gamma$
  & Agnostic real-valued learning requires margin;
    finiteness at all scales replaces a single
    finiteness condition. \\[4pt]
Agnostic PAC & Benchmark $\inf_{h \in \mathcal{H}} \risk(h)$
  & Without realizability, success is measured relative
    to the best in class. \\[4pt]
Noise model & Noise rate $\eta$
  & Labels are corrupted; sample complexity depends on
    the signal-to-noise ratio $(1-2\eta)^{-2}$. \\
\bottomrule
\end{tabular}
\end{center}

In each case, the extension is an \rel{extends\_grammar} edge in the concept
graph rather than a \rel{restricts} edge.  The distinction matters:

\begin{itemize}[leftmargin=*,itemsep=3pt]
\item A \rel{restricts} edge means the target is a special case of the source
  (no new vocabulary needed).
\item An \rel{extends\_grammar} edge means the source introduces
  \emph{new vocabulary}: concepts that cannot be expressed in the target's
  language.  The extension is irreducible.
\end{itemize}

\begin{remark}[The non-reduction principle]
\label{rmk:non-reduction}
A tempting strategy for multiclass or real-valued learning is to reduce to
binary classification (e.g., one-vs-all for multiclass, or
thresholding for real-valued).  Such reductions are computationally useful but
do not eliminate the need for new dimensions.  The VC dimension of the
reduced class does not generally equal the Natarajan, DS, pseudo-, or
fat-shattering dimension of the original class.  The reductions introduce
approximation gaps that the native dimensions avoid.  This is why the
\rel{extends\_grammar} edges exist: they mark the points where reduction to the
binary theory loses information.
\end{remark}


\begin{figure}[ht]
\centering
\begin{tikzpicture}[
  scale=0.85,
  level/.style={font=\small\sffamily, align=center},
  box/.style={
    rectangle, rounded corners=3pt, draw=defblue, thick,
    fill=defblue!5, minimum width=2.8cm, minimum height=0.8cm,
    font=\small, align=center
  },
  hubbox/.style={
    rectangle, rounded corners=3pt, draw=defblue, very thick,
    fill=defblue!12, minimum width=2.8cm, minimum height=0.8cm,
    font=\small, align=center
  },
  extarr/.style={-{Stealth[length=5pt]}, edgeorange, thick,
    decorate, decoration={snake, amplitude=1pt, segment length=6pt}}
]

% Binary core (the one machine-checked rung: fundamental_theorem, sc=0)
\node[hubbox] (bin) at (0,0) {Binary PAC\\$\VC(\mathcal{H}) < \infty$};
\node[font=\tiny\sffamily, text=defblue, anchor=north]
  at ([yshift=-2pt]bin.south) {\checkmark\,machine-checked};

% Extensions (literature characterizations; dimensions formalized as kernel defs)
\node[box] (mc) at (-4, 2.5) {Multiclass\\$\DSdim(\mathcal{H}) < \infty$};
\node[box] (rv) at (4, 2.5) {Real-valued\\$\fatshat_\gamma < \infty\;\forall\gamma$};
\node[box] (ag) at (-4, -2.5) {Agnostic\\$\VC(\mathcal{H}) < \infty$\\(same condition)};
\node[box] (no) at (4, -2.5) {Noise-tolerant\\$\VC(\mathcal{H}) < \infty$\\$+ \eta < 1/2$};

% Grammar growth arrows (each adds an irreducible primitive)
\draw[extarr] (bin) -- node[left, font=\tiny\itshape, text=edgeorange, align=center]
  {$+$ one-inclusion\\witness} (mc);
\draw[extarr] (bin) -- node[right, font=\tiny\itshape, text=edgeorange, align=center]
  {$+$ thresholds\\$+$ margins} (rv);
\draw[extarr] (bin) -- node[left, font=\tiny\itshape, text=edgeorange, align=center]
  {$+$ relative\\benchmark} (ag);
\draw[extarr] (bin) -- node[right, font=\tiny\itshape, text=edgeorange, align=center]
  {$+$ noise rate\\$\eta$} (no);

\end{tikzpicture}
\caption[Grammar growth from binary PAC learning]{Binary PAC learning sits at the
hub, characterized by finite VC dimension (\leanname{fundamental_theorem}).  Each
wavy arrow is an \rel{extends\_grammar} edge that grows the binary theory by
adding a primitive with no binary analogue: a one-inclusion witness for
multiclass learning, thresholds then a margin for real-valued learning, a relative
benchmark for agnostic learning, and a noise rate $\eta$.  The extension is
irreducible because reduction to the binary theory loses information
(\Cref{rmk:non-reduction}), proved in the multiclass case by the class with
$\Ndim=1$ and $\DSdim=\infty$ (\Cref{prop:ds-vs-nat}).}
\label{fig:grammar-growth}
\end{figure}


The grammar growth \emph{is} the extension.  In the multiclass case this
necessity is a theorem: the class with $\Ndim = 1$ and $\DSdim = \infty$ shows
that no two-function dimension can characterize learnability when $|Y| = \infty$.
For the real-valued, agnostic, and noisy settings no simpler characterization is
known, and the pattern is the same: each new primitive expresses what the existing
one cannot.


\subsection{Where Transformers Live}
\label{sec:transformer-grammar}

The extensions of this chapter are not abstractions a practitioner can avoid;
a transformer inhabits several at once.  Its next-token head is a
\emph{multiclass} hypothesis class over the vocabulary (tens of thousands of
labels, and unbounded in the limit of growing vocabularies), so it sits in
exactly the regime where the Natarajan dimension fails and the DS dimension is
the right measure.  Its logits are \emph{real-valued}, so its capacity is
governed by a scale-sensitive dimension; the certified generalization bound of
\Cref{ch:generalization} is, in this language, a fat-shattering and
covering-number control of the attention head.  And training is
\emph{agnostic}: no realizable $h^\star$ is assumed.

A softmax-attention head induces a concept class that is
well-behaved,\formal{TLT.Bridge.softAttention_wellBehaved} hence of finite VC
dimension in the binary setting.  Whether that control lifts to the multiclass
next-token head, that is, whether the DS dimension of an attention class over a
large label set is finite, is open (\Cref{open:attention-ds}).


% ================================================================
% OPEN PROBLEMS
% ================================================================

\section{Open Problems}
\label{sec:extensions-open}

Extending a characterization is where learning theory most often stumbles: a
dimension that works for binary classification need not, and the multiclass
history is the cautionary tale.  Each problem below is tagged a \emph{known
unknown} (KU), sharply posed but unresolved, or an \emph{unknown known} (UK),
where the right object or its formalization is still missing; a design-lab or
literature anchor is named where one exists.

\begin{openprob}[Tight multiclass sample complexity]
\label{open:multiclass-tight}
For finite $|Y| = k$, the best known bounds leave a $\Theta(\log k)$ gap
between the Natarajan-dimension-based upper and lower bounds on sample
complexity.  Does the DS dimension close this gap, yielding a sample
complexity of $\Theta(\DSdim / \varepsilon)$ without logarithmic factors?
\end{openprob}

\begin{openprob}[Multiclass online learning]
\label{open:multiclass-online}
The DS dimension characterizes PAC learnability in the multiclass setting.
What is the correct characterization of \emph{online} multiclass learnability?
Is there a multiclass analogue of the Littlestone dimension, and does it
relate to the DS dimension as the Littlestone dimension relates to VC?
\end{openprob}

\begin{openprob}[Formalizing the DS dimension]
\label{open:formalize-ds}
The Natarajan, pseudo-, and fat-shattering dimensions are formalized
(\leanname{NatarajanDim}, \leanname{Pseudodimension}, \leanname{FatShatteringDim}).
The DS dimension, defined by a default labeling perturbed by single-coordinate
deviations, satisfies $\Ndim \le \DSdim$ and characterizes multiclass
learnability: $\DSdim(\mathcal{H}) < \infty$ iff $\mathcal{H}$ is PAC learnable
(\cite{BrukhimCarmonDinurMoranYehudayoff2022}).  Formalize this dimension and
the characterization.
\end{openprob}

\begin{openprob}[The DS dimension of attention]
\label{open:attention-ds}
A softmax-attention head is verified well-behaved, hence finite-VC, in the
binary setting (\leanname{softAttention_wellBehaved}).  Is the DS dimension of
the multiclass next-token attention class (over a vocabulary of size $k$, and
as $k \to \infty$) finite?  A known unknown on the transformer side, and the
multiclass form of this chapter's central question.
\end{openprob}

\begin{openprob}[The fat-shattering dimension of attention]
\label{open:attention-fat}
The certified generalization bound for attention (\Cref{ch:generalization}) is
a covering-number control of real-valued logits.  Does it match the
fat-shattering dimension (\leanname{FatShatteringDim}) of the attention logit
class, and is that dimension finite at every scale $\gamma$?  A known unknown on
the transformer side.
\end{openprob}

\begin{openprob}[A realizable characterization via pseudodimension]
\label{open:pdim-realizable}
\Cref{prop:pdim-covering} bounds covering numbers, and
\Cref{thm:real-valued-char} characterizes the \emph{agnostic} case via
fat-shattering.  Is finiteness of the pseudodimension (\leanname{Pseudodimension})
exactly the characterization of \emph{realizable} real-valued learnability, and
can it be proved as a clean biconditional?  A known unknown.
\end{openprob}

\begin{openprob}[Formalizing partial-concept learnability]
\label{open:partial-formal}
The partial VC dimension characterizes partial-concept PAC learnability, with
the surprise that learnable partial classes need not extend to finite-VC total
classes~\cite{AlonHannekeHolzmanMoran2021}.  Can this be formalized, including
the no-finite-extension witness?  An unknown known: the partial-concept model
is absent from the kernel.
\end{openprob}

\begin{openprob}[List learning]
\label{open:list-learning}
A list learner outputs a short list of labels and succeeds if the true label
is among them.  What combinatorial dimension characterizes list PAC
learnability, and how does it relate to the DS
dimension~\cite{BrukhimCarmonDinurMoranYehudayoff2022}?  A known unknown.
\end{openprob}

\begin{openprob}[Agnostic proper versus improper]
\label{open:agnostic-proper}
In the realizable case, proper and improper PAC have the same sample
complexity (\Cref{sec:proper-improper}).  In the \emph{agnostic} case, is there
a class for which proper learning needs strictly more samples than improper
learning, an information-theoretic rather than merely computational separation?  A
known unknown.
\end{openprob}

\begin{openprob}[Noise-tolerant real-valued learning]
\label{open:noise-fat}
Classification noise tolerance is understood for binary VC classes
(\Cref{thm:noise-tolerant}).  What is the fat-shattering analogue: how does
label noise on real-valued targets degrade the scale-sensitive sample
complexity?  A known unknown.
\end{openprob}

\subsection*{Counterfactual exercises}

These exercises generalize the chapter's figures, each widening a concrete picture
into the chapter's argument that binary PAC theory extends only by grammar growth.

\begin{enumerate}[leftmargin=*,itemsep=3pt]
\item \textbf{The one verified rung.} The fundamental theorem, that finite VC
dimension characterizes realizable binary PAC learnability over $\mathcal{H}
\subseteq \{0,1\}^X$ (\leanname{fundamental_theorem}), is the single edge of the
concept map with a machine-checked proof.  Argue that every other characterization
in the map, multiclass, real-valued, agnostic, and proper-versus-improper, is a
literature result that grows outward from this one rung.

\item \textbf{Grammar growth as the thesis.} Read the \rel{extends\_grammar} edges
as one class.  Each marks a place where the binary language loses information and a
new primitive must be added: the pair $(f_0, f_1)$ for Natarajan
(\Cref{def:natarajan-dim}), the threshold vector then the margin $\gamma$ for the
real-valued branch (\Cref{def:pseudodim-ext}, \Cref{def:fat-shattering-ext}), and
the noise rate $\eta$ (\Cref{thm:noise-tolerant}).  Show that no reduction to the
binary theory (\Cref{rmk:non-reduction}) collapses these into VC dimension.

\item \textbf{Table and star are one object.} In Figure~\ref{fig:ds-shattering},
show that a witness table of $d$ rows, each deviating from the default $f$ at
exactly one coordinate, is the same data as a $d$-edge star out of $f$ in the
one-inclusion hypergraph $\mathrm{OIG}(\mathcal H, S)$ (\Cref{def:oig}).  Conclude
that generalizing to arbitrary $d$ means drawing a $d$-edge star.

\item \textbf{Weaker per point, larger dimension.} Using \Cref{prop:ds-vs-nat},
show every N-shattered set is DS-shattered (take $f = f_0$ and let $h_i$ realize
the basis string $e_i$), so $\Ndim(\mathcal H) \le \DSdim(\mathcal H)$.  Then
describe the Brukhim et al.\ (2022) class with $\Ndim = 1$ yet $\DSdim = \infty$,
and explain why local two-label structure fails to bound learnability when
$|Y| = \infty$.

\item \textbf{The margin is the scale primitive.} In
Figure~\ref{fig:fat-shattering}, show the symmetric band $(t_i-\gamma, t_i+\gamma)$
is the one object with no binary analogue.  Shrink it to zero width and show
$\gamma$-fat-shattering degenerates to the $P$-shattering condition of
\Cref{def:pseudodim-ext}, so $\fatshat_\gamma(\mathcal F) \to \Pdim(\mathcal F)$ as
$\gamma \to 0$.  Using \Cref{rmk:scale-family}, argue the characterizing object is
a scale family $\{\fatshat_\gamma\}_{\gamma > 0}$, one finiteness requirement per
scale.

\item \textbf{The proper-improper gap lives in the computational layer.} Trace
$3$-term DNF through Figure~\ref{fig:proper-improper}: learnable both ways
statistically with $\Theta(d/\varepsilon)$ samples, yet NP-hard properly and
polynomial-time improperly as $3$-CNF (\Cref{thm:pitt-valiant}).  Hold the columns
fixed and show that toggling only the layer turns a $\Theta(1)$ proper-to-improper
cost ratio into a superpolynomial one unless $\mathrm{RP} = \mathrm{NP}$.

\item \textbf{Is grammar growth a provable no-single-dimension invariant?} The
chapter proves grammar growth \emph{necessary} only for multiclass, via the
$\Ndim = 1$, $\DSdim = \infty$ class (\Cref{prop:ds-vs-nat}); for the real-valued,
agnostic, and noisy settings no simpler characterization is known.  Conjecture that
no single combinatorial dimension characterizes all extensions at once, a
no-free-lunch for shattering primitives, and identify what a proof would need
beyond the one multiclass instance.
\end{enumerate}

\bigskip

This chapter has traced the grammar growth required to extend binary PAC
learning theory to multiclass, real-valued, agnostic, noise-tolerant, and
computationally constrained settings.  Each extension introduces irreducible
new primitives and yields its own characterization theorem or separation
result.  The next chapter examines what lies beyond these
characterizations: the application-layer concepts, the open problems, and the
scope boundaries of the theory.
